\documentclass[11pt]{article}
\usepackage[utf8]{inputenc}
\usepackage[T1]{fontenc}
\usepackage{amsmath,amssymb}
\usepackage{graphicx}
\usepackage[margin=1in]{geometry}
\usepackage{siunitx}
\usepackage{textcomp}
\usepackage{booktabs}
\usepackage[hidelinks]{hyperref}
\usepackage{authblk}

\newcommand{\Ang}{\text{\AA}}
\newcommand{\Gmet}{\mathbf{G}}
\newcommand{\Pop}{\mathbf{P}}
\newcommand{\Uop}{\mathbf{U}}
\newcommand{\Mop}{\mathbf{M}}
\newcommand{\Bbasis}{\mathbf{B}}
\newcommand{\rootinv}{\rho}
\newcommand{\sortkey}{\mathbf{s}}
\newcommand{\Sclosure}{\mathcal{S}}

\title{\bfseries From reduced cells to lattice search:\
conservative Euclidean retrieval and exact reindexing}

\author[1]{[Author One]}
\author[1]{[Author Two]}
\affil[1]{[Affiliation], [City], [Country]}
\date{}

\begin{document}
\maketitle

\begin{center}
\begin{minipage}{0.92\textwidth}
\small\noindent\textbf{Synopsis.}
A unit cell is a basis choice, not the lattice itself. Selling reduction turns the infinite primitive-basis ambiguity into a finite closure of obtuse superbases. We use a sorted six-root-product key for continuous Euclidean retrieval, prove that its distance is a lower bound to the physically allowed relabelling-orbit distance, and certify every candidate by exact basis algebra over the full Selling-superbase closure. The resulting architecture is deliberately asymmetric: search continuously and conservatively; decide and reindex exactly.
\end{minipage}
\end{center}

\begin{abstract}
A unit cell is a choice of basis, whereas the object to be classified, compared, or reindexed is the lattice generated by that basis. Niggli reduction addressed the classical classification problem by selecting a reproducible reduced representative from infinitely many primitive bases. Modern crystallographic workloads increasingly pose a different problem: databases contain hundreds of thousands of cells, serial experiments index thousands of individual crystals, and perturbation experiments generate densely sampled lattice trajectories. These tasks require fast neighbourhood retrieval and exact basis recovery, but they do not require one representation to be simultaneously canonical, continuous, Euclidean and injective.

We formulate the computation around Selling reduction and separate three objects that serve different purposes. First, Selling reduction replaces the infinite primitive-basis ambiguity by a finite \emph{Selling-superbase closure}. For a generic Voronoi type V1 lattice there is one obtuse-superbase class up to central inversion and the 24 $S_4$ index relabellings; at the degenerate types V2--V5 additional non-isometric obtuse superbases appear, with 2, 3, 3 and 4 classes respectively in Kurlin's current classification. Second, for fast search we discard the pairing information and use the globally sorted six root products as a key $\sortkey(\Lambda)\in\mathbb R^6$. This key is continuous and Euclidean but deliberately many-to-one. By the rearrangement inequality,
\[
\|\sortkey(x)-\sortkey(y)\|_2 \leq \min_{\sigma\in G}\|x-\sigma y\|_2
\]
for any physically allowed relabelling group $G\subseteq S_6$. A radius query on a KD tree of sorted keys is therefore conservative: it cannot miss a lattice inside that radius in the corresponding orbit metric, although it may over-retrieve collisions. Third, lattice identity and reindexing are certified only after retrieval by enumerating the complete finite Selling-superbase closure, composing retained exact transformations and verifying the metric identity. The conventional-to-primitive step is rational; subsequent primitive-basis operations are unimodular integer matrices.

This architecture supports population-scale lattice geometry, crystal-by-crystal analysis and exact reindexing without asking a single coordinate system to do incompatible jobs. The full PDB can be indexed with ordinary Euclidean spatial-search machinery; serial XFEL cell variation can be analysed directly by SVD/PCA of the sorted keys; and a 200-frame monoclinic benchmark reduces an unrestricted 6960-matrix search to a small cached operator set. The conceptual rule is: \emph{search continuously and conservatively; decide exactly}. Kurlin's work supplies the root-product construction, the complete isometry classification and the finite obtuse-superbase structure. The continuity and lower-bound properties of the sorted Euclidean search key used here are elementary consequences of sorting and the rearrangement inequality; the exact operator test, not the key alone, establishes crystallographic identity.
\end{abstract}

\section{Introduction: from classification to comparison}

Crystallography has always had to distinguish a lattice from the unit cell used to describe it. If $\Bbasis$ is a primitive direct-space basis and
\begin{equation}
\Gmet = \Bbasis^{\mathsf T}\Bbasis
\end{equation}
is its metric tensor, then every unimodular integer matrix $\Mop\in \mathrm{GL}(3,\mathbb Z)$ generates another primitive basis of the same lattice,
\begin{equation}
\Bbasis' = \Bbasis\Mop, \qquad
\Gmet' = \Mop^{\mathsf T}\Gmet\Mop .
\end{equation}
The lattice is therefore not a single metric tensor but an equivalence class of metric tensors under integer change of basis.

The traditional computational response is reduction: choose one representative. Niggli reduction is the canonical crystallographic example. Its historical achievement was to turn an infinite representational ambiguity into a single reportable reduced cell. That is exactly what is wanted for classification, tabulation, comparison against reference tables, and the assignment of conventional crystallographic descriptions. The price is topological rather than numerical. A rule that selects one representative from an equivalence class must decide which representative to choose at the boundaries where two choices become equally admissible. Consequently, a smoothly changing lattice can produce a discontinuously changing reduced cell. The reduction has not made the lattice discontinuous; it has made the \emph{coordinate choice} discontinuous.

For much of crystallographic practice this is acceptable because the dominant question is classificatory: what lattice is this, and what conventional cell should be reported? Modern computation increasingly asks a different question: how close are two lattices? Structural archives contain hundreds of thousands of reported cells. Serial crystallography produces thousands to millions of indexing solutions. Thermal, pressure, hydration, ligand-binding and operando experiments produce trajectories of related lattices. In these settings a discontinuous representative is awkward because the computational task is not merely to attach one label to each lattice but to compare large populations, construct neighbourhoods, follow continuous paths, and recover basis transformations only when required.

The central argument of this paper is that modern lattice computation should not force one representation to perform three incompatible tasks. The complete Kurlin root form is designed to distinguish isometry classes. The continuous Euclidean object that is most convenient for large-scale search is instead a deliberately coarser sorted root-product key. Exact crystallographic identity and reindexing are then decided in a third representation: the finite Selling-superbase closure with retained basis transformations.

Selling reduction is the pivot. It replaces infinitely many primitive bases by finitely many obtuse superbases. For a generic lattice this finite ambiguity is generated by the 24 relabellings of one obtuse superbase (up to central inversion). At degenerate Voronoi types the closure becomes richer rather than smaller: additional obtuse superbases not related by $S_4$ relabelling must also be included. Kurlin's current classification describes these classes explicitly. This finite closure is the basis-aware object on which exact operator recovery is performed.

For search we take a different projection of the same root products. Let $x\in\mathbb R_{\ge0}^6$ be the six root products of an obtuse superbase and let $\sortkey(x)=\operatorname{sort}(x)$ be their nondecreasing ordering. This key forgets which tetrahedral edge carried which root product. It is therefore not a complete lattice invariant: distinct lattices can collide. That loss of information is useful rather than dangerous because the key is a Euclidean lower bound to every physically allowed relabelling-orbit distance. Radius search in key space can over-retrieve but cannot under-retrieve relative to that orbit metric. Exact basis algebra then certifies the survivors.

A lattice search key alone does not solve crystallographic reindexing. Reindexing requires a basis-level map. We therefore keep the exact transformations used to reach the primitive and Selling-reduced descriptions. The conventional-to-primitive step can require a rational matrix whose denominators are fixed by centering; once the primitive lattice has been reached, reduction and reindexing proceed with unimodular integer matrices. A candidate pair retrieved in Euclidean key space is lifted through the complete finite Selling-superbase closure and accepted only if an exact integer change of basis satisfies the metric relation. This gives the paper's computational division of labour:
\begin{quote}
\textbf{Search conservatively in $\mathbb R^6$; certify and operate in $\mathrm{GL}(3,\mathbb Z)$.}
\end{quote}
Continuous Euclidean keys are used for neighbourhood retrieval, exploratory graphs and linear population analysis. Exact rational/integer algebra is used when lattice identity, atoms, Miller indices, symmetry operations or indexing solutions must be transformed.

The mathematical attribution is correspondingly narrow. Kurlin establishes the complete 3D isometry classification through root forms, describes the type-dependent families of obtuse superbases, and bounds the change of individual root products under perturbation. His current manuscript explicitly postpones the construction of a metric on the 3D root-invariant space. We therefore do not identify the complete ordered root form with the Euclidean search metric used here. The sorted Euclidean key, its continuity at ties, and its conservative lower-bound property are the computational construction used in this work; exact identity is established independently by the operator test.

\section{Three objects: Selling closure, Euclidean key, exact operator}

\subsection{One controlled rational step, then exact integer lattice algebra}

A reported crystallographic cell is not always primitive. Converting a centered conventional cell to a primitive basis can therefore require a rational change-of-basis matrix. This is the only point in the workflow at which fractions are structurally unavoidable. The denominators are small and determined by the centering relation. Once a primitive basis has been obtained, changes between primitive bases of the same lattice are unimodular integer transformations.

This separation is useful computationally. The conventional-to-primitive map is retained exactly as a rational matrix. Selling reduction then accumulates only integer elementary transformations. All subsequent reindexing operators between primitive bases are exact integer matrices. Compositions therefore do not accumulate floating-point drift, and the same operators can be applied consistently to reciprocal indices, direct-space fractional coordinates, orientation matrices and symmetry operators.

\subsection{Selling reduction makes the equivalence problem finite}

Given primitive basis vectors $v_1,v_2,v_3$, introduce
\begin{equation}
v_0 = -(v_1+v_2+v_3),
\end{equation}
so that the four vectors form a superbase. Define the six Selling conorms
\begin{equation}
p_{ij} = -v_i\cdot v_j, \qquad 0\le i<j\le 3.
\end{equation}
In the Andrews--Bernstein--Sauter $S^6$ convention the Selling scalars occupy the non-positive orthant, $s_{ij}=-p_{ij}$~\cite{abs2019}; the search key used below is therefore $\operatorname{sort}(\sqrt{-s})=\operatorname{sort}(\sqrt{p})$. All six scalars are homogeneous (dot products of the same dimension); taking square roots~\cite{kurlin2026} places them in length units, which makes the volume normalisation $\sortkey/V^{1/3}$ dimensionally clean. A superbase is obtuse when all six conorms are non-negative. Selling reduction reaches such a superbase by a finite sequence of integer operations: the practical Delone move (negate the positive scalar product, adjust the four adjacent products, and swap two vectors) is the crystallographic algorithm of Andrews and Bernstein~\cite{ab1988,ab2014,abs2019}, while Kurlin's Lemma~A.1 supplies the termination argument that the sum of the six conorms decreases at each step~\cite{kurlin2026}.

The conceptual gain is that the infinite primitive-basis ambiguity becomes finite, but the finite object is more subtle than one $S_4$ orbit at the symmetry boundaries. Classical crystallography already records that the six reduced Selling scalars are unique as a list (a multiset), while their arrangement is unique only up to the 24 index relabellings of the four superbase vectors~\cite{bravais1850,delone1932,abs2019}. Kurlin's Lemmas 4.1--4.5 refine the arrangement side: for generic type V1, a lattice has one pair of centrally symmetric obtuse superbases and the associated coforms are related by those 24 $S_4$ permutations, but at the higher-symmetry/degenerate types there are additional \emph{non-isometric obtuse-superbase classes}: 2 for V2, 3 for V3, 3 for V4 and 4 (= $1+3$) for V5 in the April 2026 classification. For a primitive orthorhombic lattice, for example, the V5 closure contains 32 obtuse superbases split into four isometry classes. In short: Andrews--Bernstein--Sauter state the multiset uniqueness; Kurlin gives the full arrangement closure.

We therefore distinguish an $S_4$ \emph{relabeling orbit} from the full \emph{Selling-superbase closure} $\Sclosure(\Lambda)$. The latter contains every inequivalent obtuse superbase required by the Voronoi type together with its permitted index relabellings (and central inversion where relevant), deduplicated at the level needed by the application. This is the finite basis-aware object used by the exact reindexing stage. Treating the closure as ``at most 24'' is correct only for generic V1 and is insufficient for orthorhombic, tetragonal, hexagonal and cubic loci that are common in crystallographic databases.

Selling reduction therefore need not provide one globally preferred representative. Its computational role is to turn an infinite basis problem into a finite closure while retaining exact integer provenance for every branch that may later be required.

\subsection{Complete root form and continuous Euclidean search key are different objects}

For an obtuse superbase, define the root products
\begin{equation}
r_{ij}=\sqrt{p_{ij}}\ge 0.
\end{equation}
Kurlin arranges these quantities in a type-dependent $2\times3$ root form that retains the pairing of opposite tetrahedral edges. His classification uses that structured form, together with the finite alternatives described above, to distinguish 3D lattice isometry classes. We rely on that result for completeness of the classification, not for the Euclidean metric used in our database search.

For retrieval we deliberately use a coarser object. Write the six root products as a vector $x=(r_{01},r_{02},r_{03},r_{12},r_{13},r_{23})$ in any order and define
\begin{equation}
\sortkey(x)=\operatorname{sort}(x)\in\mathbb R_{\ge0}^6,
\end{equation}
where the components are arranged nondecreasingly. The key is invariant to all permutations of the six entries and is therefore more quotienting than the physical $S_4$ action on tetrahedral edges. It is not injective in V1, V2 and V4: the forgotten pairing information produces finite collisions (generically 30-to-1, 15-to-1 and 4-to-1, respectively, for the corresponding pairing patterns). It is injective in V3 and V5. These multiplicities are a feature of the deliberately coarse key, not a failure of the exact certification stage. We do not use equality of $\sortkey$ as a proof of lattice identity.

The benefit is that $\sortkey$ is continuous and Euclidean. Kurlin's Lemma~6.6 bounds changes of individual root products under perturbation of a fixed obtuse superbase by $\sqrt{2\ell\delta}$, which is \emph{H\"older-$\tfrac12$} rather than Lipschitz at vanishing conorms: angular noise that is Gaussian in conorm space ($p_{ij}$ is linear in the metric) is amplified by $\sqrt{\,\cdot\,}$ exactly on the zero slots of high-symmetry cells. On the CXIDB~83 hexagonal cell $(41.8,41.8,233,90,90,120)$, a $0.01^\circ$ angular perturbation already shifts the sorted $\sqrt{\,}$ key by a median $\sim1.4\,\Ang$ --- more than ten times the reported $c$-axis MAD of $0.13\,\Ang$ --- whereas the same noise moves a generic triclinic key by only $\sim0.3\,\Ang$ at $0.10^\circ$. Sorting does not introduce a tie discontinuity: exchanging the order of two equal or crossing components leaves the sorted vector continuous. Continuity also survives a change of Selling branch at a Voronoi-type boundary. Kurlin's Lemmas 4.2--4.5 show that the alternative obtuse superbases in a degenerate closure have coforms related by transpositions or permutations of their entries, and therefore all members of the closure have the same sorted root key. A perturbation crossing the boundary may reduce to a superbase near a different closure member, but whichever branch is selected its root products approach the same sorted limiting key; combined with the root-product perturbation bound, the lattice-level key therefore varies continuously through the boundary, albeit with the H\"older-$\tfrac12$ rate at symmetry strata. More strongly, the sorted distance is a certified lower bound on any physically allowed relabelling-orbit distance.

Any monotone per-slot map $f$ that depends only on the conorm $p$ and on lattice invariants (not on a particular superbase's edge lengths) preserves the properties we rely on when applied before sorting: relabelling permutes slots so sorting still yields one key per lattice, and the rearrangement lower bound holds for $\operatorname{sort}(f(x))$ versus $\operatorname{sort}(f(y))$. Three practical stabilisations of $\sqrt{p}$ are therefore available. The noise floor must itself be closure-invariant: the natural choice is the global scale $s=\sigma_\theta\,T$ with $T=\sum_{ij}p_{ij}=\tfrac12\sum_i|v_i|^2$, which is unchanged by Selling flips (Kurlin Lemma~A.1 with $\varepsilon=0$ permutes the conorms). Per-pair floors $|v_i||v_j|\sigma_\theta$ are \emph{not} invariant across the closure and must not be used.
\begin{itemize}
\item floored root $\sqrt{p+s}-\sqrt{s}$ (Wiener-style; Lipschitz near zero; $\sim5\times$ noise reduction; length units);
\item soft threshold $\sqrt{\max(p-\kappa s,0)}$ (shrinkage onto the symmetry stratum; $\sim15\times$ reduction; deliberate blindness below $\kappa\sigma_p$);
\item linear conorm key $\operatorname{sort}(p)$ or $\operatorname{sort}(p/\sqrt{T})$ (Lipschitz, Gaussian-noise preserving; natural for per-frame PCA).
\end{itemize}
\begin{table}[h]
\centering
\caption{Median sorted-key shift (\Ang) under angular noise on the $P6_3$ cell $(41.8,41.8,233,90,90,120)$, comparing plain $\sqrt{p}$ with the three stabilisations (invariant floor $s=\sigma_\theta T$; soft threshold $\kappa=2$; linear uses $p/\sqrt{T}$).}
\begin{tabular}{ccccc}
\toprule
$\sigma_\theta$ & $\sqrt{p}$ & $\sqrt{p+s}-\sqrt{s}$ & $\sqrt{\max(p-2s,0)}$ & $p/\sqrt{T}$ \\
\midrule
$0.01^\circ$ & 1.58 & 0.28 & 0.03 & 0.01 \\
$0.05^\circ$ & 3.12 & 0.64 & 0.17 & 0.06 \\
$0.10^\circ$ & 4.63 & 0.68 & 0.35 & 0.10 \\
\bottomrule
\end{tabular}
\end{table}
A stabilised key is a different metric from Kurlin's: choosing $s$ sets the resolution at which ``symmetric'' and ``nearly symmetric'' stop being distinguished, and the conservative-retrieval lemma then holds in the stabilised metric. The archival search key in this paper remains the default $\sqrt{p}$; for noisy serial frames we prefer the sorted-conorm (linear) key.

\paragraph{Lemma (conservative sorted-key distance).}
Let $x,y\in\mathbb R^6$ and let $G\subseteq S_6$ be any allowed set of permutations of the six root products. Then
\begin{equation}
\label{eq:lowerbound}
\|\sortkey(x)-\sortkey(y)\|_2
=\min_{\sigma\in S_6}\|x-\sigma y\|_2
\leq \min_{\sigma\in G}\|x-\sigma y\|_2.
\end{equation}
The equality is the rearrangement inequality applied to the squared Euclidean distance, and shows that this is the tightest lower bound obtainable when only the unpaired six-entry multisets are retained. Because the physically realizable superbase relabellings form only a subset of all six-entry permutations, forgetting the edge pairing can only decrease the distance.

Equation~\eqref{eq:lowerbound} gives a useful one-sided guarantee. If the true relabelling-orbit distance between two root-product lists is at most $r$, then their sorted-key distance is also at most $r$. A KD-tree \emph{radius} query in sorted-key space therefore cannot miss such a candidate; it may only return additional collisions. Those false candidates are removed by the exact operator stage. Ordinary fixed-$k$ nearest-neighbour queries are useful exploratory tools but do not by themselves preserve the ranking of the true orbit metric; certified nearest-neighbour retrieval requires radius expansion or over-retrieval followed by exact certification.

At degenerate Voronoi types the full Selling-superbase closure contains several inequivalent obtuse superbases, but this does \emph{not} multiply the search keys. The ABS/Bravais multiset uniqueness already implies that every obtuse presentation of a lattice carries the same unordered list of six scalars; Kurlin's Lemmas 4.2--4.5 relate the coforms of the alternative classes by transpositions or permutations of entries, so every obtuse superbase of a given lattice has the same sorted multiset of root products. There is therefore exactly \emph{one sorted Euclidean key per lattice}, for V1 through V5. The expanded closure is needed only after retrieval, when an exact basis correspondence and operator are required.

The complementary seven vonorms of an obtuse superbase --- the squared lengths of the four superbase vectors and of the three opposite-edge partial sums --- are the $D_7$ descriptor of Andrews--Bernstein--Sauter~\cite{abs2019} and Kurlin's voform~\cite{kurlin2026}. Kurlin's Example~6.5 shows that $D_7$ alone is incomplete: two non-isometric lattices can share identical vonorm multisets while carrying distinct conorm multisets, which the sorted six-root key separates. The converse also holds. On a generic V1 fibre the sorted six-key merges 30 lattices that differ by tetrahedral edge pairing; in 300 random such fibres every member had a distinct sorted-$\sqrt{D_7}$ key. The concatenated Euclidean key
\[
\operatorname{sort}(r)\;\big\|\;\operatorname{sort}(\sqrt{\text{vonorms}})
\in\mathbb R^{13}
\]
is therefore still a pure sort --- continuous, with the lower-bound property of equation~\eqref{eq:lowerbound} holding blockwise --- and collapses the generic over-retrieval multiplicity to essentially one at negligible cost. This is an empirical strengthening rather than a proof (ties and degenerate types need separate treatment); the primary archival search key in this paper remains the six-dimensional sorted root product.

For shape-only retrieval we remove overall scale from each key, for example by
\begin{equation}
\tilde{\sortkey}(\Lambda)=\frac{\sortkey(\Lambda)}{V(\Lambda)^{1/3}}.
\end{equation}
The resulting vector remains an ordinary point in Euclidean space. KD trees, radius graphs, PCA/SVD and other vector algorithms can therefore be used as fast retrieval and exploratory layers without conflating the search key with a complete lattice invariant.

\subsection{Exact certification and lift through the complete Selling closure}

A sorted key deliberately removes basis labels and, generically, some lattice-identifying pairing information. It is therefore never the final identity test. During Selling reduction we retain the integer coordinates of every superbase branch needed for $\Sclosure(\Lambda)$ with respect to the input primitive basis. For a retrieved pair $A,B$, let $\Uop(A)$ and $\Uop(B)$ contain three independent vectors of a candidate pair of Selling superbases, expressed in their respective original primitive bases. The candidate change of basis is
\begin{equation}
\Pop = \Uop(A)\Uop(B)^{-1}.
\end{equation}
It is accepted only if
\begin{equation}
\label{eq:metricverify}
\Pop^{\mathsf T}\Gmet(A)\Pop = \Gmet(B)
\end{equation}
for exact metrics, or by an explicitly reported residual for noisy measured metrics. Because $\Uop(A)$ and $\Uop(B)$ are unimodular integer matrices after the primitive conversion, $\Pop$ is an exact integer operator. Floating-point arithmetic is used to assess measured metrics, not to invent or round the operator.

The enumeration must run over the complete finite Selling-superbase closure of both lattices. For generic V1 this reduces to the familiar $S_4$ relabellings (plus central inversion as appropriate). For V2--V5 it must also include the alternative non-isometric obtuse superbases classified by Kurlin. This point is operationally important: restricting the lift to 24 relabellings of a single reduced superbase is not closure-complete for many high-symmetry crystallographic lattices.

If the lattice has non-trivial automorphisms, the correct result is a set or coset of symmetry-related operators rather than a unique matrix. In serial crystallography these residual cosets are the indexing branches that geometry alone cannot distinguish and that intensity correlations may subsequently resolve.

The architecture is therefore a filter-and-certify pipeline. The sorted key answers the permissive question, \emph{which lattices could be within this radius?}; the Selling closure and equation~\eqref{eq:metricverify} answer the exact question, \emph{are these the same lattice and, if so, which basis operator connects them?}

\section{Results}

\subsection{A representative can jump while the lattice does not}

We first isolate the reason a single reduced cell is an awkward object for continuous comparison. A synthetic monoclinic series was generated with $a=120\,\Ang$, $b=189.1\,\Ang$, $\beta=91.2^\circ$ and $c$ varying continuously from $150$ to $100\,\Ang$. The path passes through the pseudo-symmetric and then symmetric region around $a=c$.

Across such a trajectory the underlying metric changes continuously. A reduced-cell prescription, however, can change which basis it reports when the trajectory crosses a reduction boundary. The discontinuity is potentially large in the coordinates of the selected representative even when the perturbation in the lattice is arbitrarily small. In the boundary-straddling test used here, the distance between raw reduced $S^6$ representatives changed from 0.08 to 83.7 for the same physical perturbation on opposite sides of the boundary, whereas the sorted-root key follows the perturbation continuously. That jump is precisely the arrangement ambiguity that Andrews--Bernstein--Sauter acknowledge: a single $S^6$ representative is one of 24 (or more) arrangements, so raw $S^6$ distance is not itself a lattice distance~\cite{abs2019}. Their companion treatment restores a lattice metric by minimising over reflections and boundary transforms. The sorted key removes that orbit minimisation by construction, at the cost of injectivity; they keep injectivity at the cost of an orbit search. That is a genuine trade-off, not a defect in $S^6$. The previously reported numerical value 0.0010 must be rechecked against the exact sorted-key implementation before submission and is therefore not used as evidence here.

The point of this experiment is not to re-prove the continuity theorem. It makes the crystallographic consequence visible: \emph{the discontinuity belongs to the representative, not to the lattice}. The finite Selling-superbase closure contains the competing presentations on both sides of the transition, and the sorted root key provides a continuous retrieval coordinate while the Selling closure retains the basis alternatives needed for exact certification.

The same example also clarifies why symmetry loci should not be treated as pathological exceptions to continuous retrieval. The point $a=c$ is precisely where an axis exchange becomes an exact lattice symmetry. In a trajectory description, such a point is a junction at which the available basis correspondences change. A useful lattice coordinate should allow the path to pass through that junction continuously, while the exact-basis layer should expose the additional operators that become admissible there.

% FIGURE 1 SUGGESTION:
% Same underlying monoclinic lattice trajectory shown in three rows:
% (a) selected Niggli/Selling representative with a visible flip,
% (b) complete finite Selling-superbase closure around the boundary,
% (c) continuous sorted-key trajectory.
% Add the exact integer flip operator at the crossing.

\subsection{Sorted root keys turn archival lattice retrieval into ordinary computational geometry}

The practical importance of the root representation becomes clearest at archive scale. We computed six-dimensional root search keys for 206\,214 crystallographic cells from the PDB. Because the keys are ordinary Euclidean vectors, the database can be indexed with \texttt{scipy.spatial.cKDTree}. The previous implementation reported an index-build time of \SI{0.28}{\second} and median key-space $k=10$ query time of \SI{0.04}{\milli\second}; a 3000-cell benchmark reported \SI{0.04}{\second} build time and \SI{0.02}{\milli\second} median key-space query time. Before submission these timings will be rerun explicitly on the sorted key defined in Section~2.3. For scale, the SAUC search of alternate unit cells reported 500-nearest-neighbour queries over roughly half a million PDB+CSD cells in about \SI{4}{\second} real time with a Selling embedding versus about \SI{8}{\second} with Niggli~\cite{mcgill2014}; those times search an orbit-aware space, whereas the millisecond figures above search the deliberately lossy sorted key and therefore require the subsequent exact certification stage.

The significance is not the specific KD-tree implementation. The significance is that the high-volume retrieval stage can be made an ordinary Euclidean search. The KD tree does not prove lattice identity and a fixed-$k$ query is not a certified ranking of the true orbit metric. For a certified threshold query we instead use the lower-bound property of equation~\eqref{eq:lowerbound}: a radius search returns a conservative candidate superset, after which the complete Selling closure and exact operator test remove collisions. There is one sorted key per lattice at every Voronoi type; the additional V2--V5 closure branches are never inserted into the index and are opened only for exact certification after retrieval.

For shape visualization we use the volume-normalised key
\[
\tilde{\sortkey}=\sortkey/V^{1/3}.
\]
This places cells in a common Euclidean \emph{sorted-root key space} after removal of overall scale. We used the complete PDB population to build a metric-neighbourhood graph and a two-dimensional visual embedding. It is important to state exactly what is being embedded. The input is not Kurlin's complete pairing-preserving root invariant and the resulting map is not claimed to be an embedding of the complete 3D lattice-isometry space. It is an embedding of the Euclidean pseudometric induced by the deliberately coarser sorted key.

That distinction is operationally useful rather than merely cautionary. Sorting is what gives the representation the continuity and ordinary Euclidean geometry needed by KD trees, neighbourhood graphs, SVD/PCA and nonlinear visualization. The price is finite non-injectivity: two genuinely different lattices can have the same sorted key because the tetrahedral pairing information has been discarded. No downstream embedding can recover information that was removed at this stage. Accordingly, a collision in sorted-key space is treated exactly as it is in database search: as an unresolved equivalence that requires the complete Selling closure and exact operator test if crystallographic identity matters.

The embedding therefore has a deliberately narrower scientific interpretation. It is not evidence that nearby lattices are nearby; that relationship is imposed by construction through the sorted-key distance. Nor can a branch, junction or apparent connection by itself establish the topology of the complete lattice-isometry space. The scientifically interesting information is how the \emph{observed crystallographic population} occupies this continuous Euclidean quotient, especially when structure in that population is explained by labels or experimental variables that were not used to construct the key.

When the embedding is coloured after construction by crystal system and symmetry information, coherent regions and junctions appear. Tetragonal populations form branches that approach cubic symmetry at a restricted locus; orthorhombic populations occupy neighbouring and connecting regions; trigonal and hexagonal regions appear at other boundaries; and monoclinic cells populate a broader portion of the space. These labels were not used to construct the sorted-key distance. Thus the visualization provides a population-level view in which familiar discrete crystallographic classes appear as structure within the continuous sorted-root representation. We deliberately avoid interpreting these structures as a proof of the global topology of lattice-isometry space until the effect of sorted-key collisions on the local graph has been quantified.

A direct collision audit provides that check. For every exact or near collision of sorted keys, and more generally for edges of the $k$-nearest-neighbour graph, the candidate pair can be passed through the complete Selling-superbase closure and exact operator machinery. This distinguishes three cases: identical lattices represented by different bases; distinct lattices that collide under sorting; and ordinary nearby lattices whose local ordering is unaffected by the lost pairing information. Reporting the fraction of local graph edges preserved after closure-aware certification therefore measures how faithfully the sorted-key geometry represents the \emph{occupied} PDB population, without confusing empirical faithfulness with global injectivity. This audit is a required validation of the final embedding analysis.

To quantify how many independent shape directions are occupied locally, we compute a mean-centred SVD of the six-dimensional normalised sorted keys within neighbourhoods of the embedding. Using the entropy effective rank
\begin{equation}
r_{\rm eff}=\exp\left(-\sum_i p_i\log p_i\right),\qquad p_i=\frac{\sigma_i}{\sum_j\sigma_j},
\end{equation}
we observe local ranks from approximately 1 to 5.1, with median 4.12 over occupied regions. Low-rank filaments correspond to near one-parameter families; higher-rank regions sample a fuller set of lattice-shape degrees of freedom. This local analysis is performed in the original sorted-key space; the two-dimensional embedding is only a visual map on which to display it.

% FIGURE 2 SUGGESTION:
% Full-PDB root-shape embedding coloured by crystal system / Bravais class.
% Insets: cubic pinch, tetragonal wide-vs-long branches, orthorhombic bridge,
% trigonal/hexagonal boundary. Second panel: local SVD effective rank.

\subsection{Archival protein families occupy continuous lattice trajectories}

To examine the local structure of an experimentally familiar family, we collected every hen egg-white lysozyme entry in the PDB (UniProt P00698): 1372 entries, of which 1361 have a crystallographic unit cell. The set spans six space groups and four crystal systems, including 1152 tetragonal $P4_32_12$ cells, 100 orthorhombic $P2_12_12_1$ cells, 59 monoclinic $P2_1$ cells, 41 triclinic $P1$ cells and smaller trigonal/hexagonal populations.

The value of the invariant here is not that the dominant tetragonal cells become close---that follows directly from their similar metrics. Rather, the invariant supplies one setting-independent comparison space in which all of the reported cells can be handled together before any crystal-system-specific analysis is chosen. Within the dominant tetragonal population, 1113 native-volume cells span a 26\% range in unit-cell volume. The leading intrinsic coordinate of their sorted-key distance graph correlates with the cell-shrinkage coordinate at $|r|=0.92$. Along that coordinate, $a$ varies smoothly from approximately $76.8$ to $80.0\,\Ang$, while $c$ remains comparatively flat and then shifts between preferred values near 37 and $38\,\Ang$.

For this high-symmetry single-setting subset, the same trend can of course be described directly in $(a,c)$. The point is therefore not that the sorted root key reveals a hidden two-parameter tetragonal fact. The point is that the family was isolated from a heterogeneous archival census using one setting-independent retrieval distance, after which conventional physical interpretation can proceed in the coordinates most appropriate to that subset. The archive becomes a set of occupied trajectories and regions in lattice space rather than a collection of unrelated reported cells.

\subsection{Serial crystallography: linear modes of variation in root space}

A stronger test of the representation is provided by serial data, where every diffraction pattern may carry its own indexed cell. Merged depositions generally preserve only the final consensus cell, hiding the distribution on which indexing, outlier rejection and merging decisions were made. We parsed the CrystFEL stream from a $\beta$-lactamase serial XFEL data set (CXIDB 83): 14\,445 diffraction patterns, 12\,474 indexed crystals, space group $P6_3$, with a typical cell near $41.8\times 41.8\times233\,\Ang$.

The 12\,474 indexed cells form a compact cloud whose variation can be analysed directly in six dimensions. Because this is a high-symmetry ($P6_3$) stream with unconstrained CrystFEL angles near the vanishing-conorm strata, we key each frame by its \emph{sorted conorms} rather than by $\sqrt{p}$: the linear key is Lipschitz in the metric and does not amplify angular noise through the H\"older-$\tfrac12$ cusp discussed in Section~2.3. For a single experiment we prefer an explicit linear decomposition to a nonlinear embedding: after centering the sorted-conorm keys, SVD/PCA identifies the dominant directions of lattice variation and their associated loadings.

The long $c$ axis has a sharp central distribution, with median $233.3\,\Ang$ and median absolute deviation $0.13\,\Ang$, but a heavy tail extends to approximately $245\,\Ang$. Because the lattice key is blind to crystal orientation, orientation metadata can be projected onto the lattice-variation coordinates as an independent diagnostic. Crystals for which the incident beam lies within $10^\circ$ of $c^*$ show approximately $1.53\,\Ang$ scatter in the inferred long axis, compared with roughly $0.5\,\Ang$ for crystals oriented near the $ab$ plane. Thus the long axis is measured about three times less precisely when the beam looks down it. (These association numbers should be rechecked under the sorted-conorm key before submission; the qualitative orientation diagnostic remains the intended test.)

This association is not built into the lattice metric. It is an external experimental variable explaining part of the observed population structure. The result therefore separates two effects that a single merged cell cannot: genuine lattice variation and orientation-dependent measurement uncertainty. The same analysis could be extended to pulse energy, detector geometry, time delay, temperature or other per-pattern metadata without changing the lattice representation.

% FIGURE 3 SUGGESTION:
% XFEL sorted-key PCA/SVD scores. Panel A: PC1/PC2 scores coloured by c.
% Panel B: PC loadings on the six root components.
% Panel C: c-axis scatter versus beam angle to c*.

\subsection{Reindexing by identifying the lattice first}

The same construction gives a direct reindexing strategy. A conventional approach starts with two basis descriptions and asks which small integer matrix maps one to the other, often after reduction and symmetry analysis. That formulation mixes lattice identification with basis recovery. Here they are separated.

For each input cell we first convert to a primitive basis and retain that exact rational transformation. We then Selling-reduce while accumulating unimodular integer operations. The sorted root key is used only to retrieve candidate lattice matches. Once a candidate has been returned, identity and the operator are determined by the complete finite Selling-superbase closure, including alternative obtuse superbases at V2--V5 as well as $S_4$ relabellings. Their retained integer coordinates yield exact candidate reindexing matrices, which are verified against the metric. Symmetry-related candidates are kept as a coset for subsequent intensity-based discrimination.

The finite nature of the Selling-superbase closure makes this particularly attractive for repeated measurements around a reference lattice. The admissible operator set can be computed once for the reference and cached. A new frame is then assigned relative to that finite set rather than starting a new unrestricted matrix search.

On 200 perturbed monoclinic frames spanning both sides of a reduction flip, a brute-force reference implementation tested 6960 small unimodular matrices per frame and required \SI{23}{\milli\second} per frame. The root-first workflow identified the lattice, lifted each frame through its relevant Selling-superbase closure, and reduced the admissible result to a cached two-operator reference coset requiring \SI{0.07}{\milli\second} per frame, approximately a 330-fold reduction in per-frame cost (approximately 180-fold amortised in the present implementation).

The conceptual gain is more important than the benchmark. Reindexing is no longer posed as ``search the basis group until something fits.'' It is posed as ``identify the lattice, then resolve the basis within a finite known closure.'' This is exactly the division of labour implied by the representation.

\subsection{Reduction flips are recovered as exact operators, not tolerated as pseudo-symmetry}

Near pseudo-symmetric boundaries, tolerance-gated symmetry enumeration can drop a basis transformation that is required only because the reduction selected a different representative. The distinction is important: the surface at which a reduction flips and the locus at which the corresponding basis exchange becomes an exact metric symmetry need not coincide.

In a controlled monoclinic family, exchanging the $a$ and $c$ axes is an exact integer relation between two settings of the same lattice for every value of $|a-c|$, but it is an exact metric symmetry only at $a=c$. As $|a-c|$ increases, the Le Page angular residual of the exchange grows. With a $3^\circ$ tolerance, the operator is admitted up to approximately $|a-c|=2.85\,\Ang$ and then excluded, even though the two descriptions remain exactly related by the same integer change of basis. Raising the tolerance merely admits additional unrelated pseudo-symmetries.

The Selling/root workflow treats the problem differently. The two settings are admitted by the conservative key search, and the complete finite superbase correspondence recovers the exact integer transformation independently of whether that transformation is close to an exact metric automorphism. Symmetry tolerance is therefore not asked to solve a reduction bookkeeping problem.

\begin{table}[t]
\centering
\caption{Monoclinic basis exchange under increasing metric distortion. The two settings describe the same lattice throughout. A $3^\circ$ pseudo-symmetry gate eventually drops the exchange, whereas complete Selling-closure recovery continues to return the exact integer operator.}
\begin{tabular}{cccccc}
\toprule
$|a-c|$ (\Ang) & flip $\delta$ ($^\circ$) & sorted-key dist. & $\delta$-gate & exact op. & residual \\
\midrule
0.0  & 0.00  & $0$ & keeps & yes & $0$ \\
1.0  & 1.04  & $7\times10^{-15}$ & keeps & yes & $0$ \\
2.0  & 2.09  & $2\times10^{-14}$ & keeps & yes & $0$ \\
2.85 & 2.98  & $7\times10^{-15}$ & keeps & yes & $0$ \\
4.0  & 4.18  & $1\times10^{-14}$ & drops & yes & $0$ \\
5.0  & 5.22  & $2\times10^{-14}$ & drops & yes & $0$ \\
10.0 & 10.41 & $7\times10^{-15}$ & drops & yes & $0$ \\
\bottomrule
\end{tabular}
\end{table}

\section{Discussion}

\subsection{The historical problem changed}

Reduction methods are sometimes compared as though they were competing answers to one timeless question. The more useful perspective is that different constructions solve different computational problems.

The classical crystallographic need was classification and standardisation. From infinitely many descriptions of one lattice, choose a reproducible representative that can be tabulated and related to Bravais type and conventional setting. Niggli reduction is a powerful answer to that problem. Its seams are not evidence that it failed; they are the unavoidable consequence of selecting one representative from an equivalence class.

The modern comparison problem is different. Given two measured lattices, a database of hundreds of thousands of lattices, or a trajectory of thousands of related measurements, we want to know which lattice each observation represents, which observations are close, how a population is structured, and only then which basis transformation is required for a particular operation. In that setting the finite Selling-superbase closure is a more natural basis-aware intermediate object than a single selected cell, while a coarser sorted Euclidean key is a more natural object for high-volume retrieval.

Selling reduction makes the representational ambiguity finite, but the exact closure is type-dependent. Kurlin's work tells us which obtuse-superbase classes must be retained for completeness. We then intentionally project their root products to a coarser sorted Euclidean key for search. The resulting order of computation is:
\begin{enumerate}
\item remove conventional centering and enter primitive lattice space;
\item construct the finite Selling-superbase closure appropriate to the Voronoi type;
\item use exactly one sorted root key per lattice for conservative Euclidean radius retrieval;
\item certify returned candidates by exact metric/operator tests over the full closure;
\item propagate the recovered rational/integer basis operator to the crystallographic data that require it.
\end{enumerate}
This is not a replacement for Niggli reduction. It is a different computational route for a different primary task.

\subsection{A Euclidean lower-bound key is enough for fast certified retrieval}

The useful computational fact is not that the complete lattice-isometry space has been equipped here with an ordinary $L^2$ metric; it has not. The useful fact is that a deliberately lossy key has a Euclidean distance that lower-bounds the physically allowed relabelling-orbit distance. That is sufficient for a certified filter. Standard KD-tree radius search can reject the overwhelming majority of a large archive without risking false negatives within the requested orbit-distance threshold, while the exact crystallographic stage handles the small over-retrieved set.

This distinction matters for how results are stated. Key-space $k$-nearest-neighbour queries are fast exploratory operations but need not preserve the true orbit ranking because collisions can crowd the list. Radius queries, or an adaptive over-retrieval scheme with subsequent certification, carry the one-sided guarantee. PCA/SVD, local-dimensionality calculations and embeddings on the sorted keys are likewise descriptive tools for the Euclidean retrieval representation, not proofs about a complete metric on lattice-isometry space. This is not a weakness of the embedding algorithm: any two lattices identified by the sorted key are already indistinguishable before dimensionality reduction. The appropriate validation is therefore performed upstream, by auditing sorted-key collisions and local graph edges against closure-complete exact comparisons.

\subsection{Continuous lattice space does not make discrete crystallography disappear}

A continuous representation should not be mistaken for a claim that Bravais classes or crystal systems are arbitrary. On the contrary, symmetry is expected to appear as special structure within lattice space: higher-symmetry lattices occupy lower-dimensional loci at which additional metric equalities hold and additional automorphisms become available. The observed organisation of the PDB population is therefore naturally read as a continuous geometric object stratified by familiar discrete symmetry classes.

This viewpoint is useful because experimentally observed lattice changes need not remain inside one class. A trajectory may approach or pass through a higher-symmetry locus, bifurcate into lower-symmetry regions, or sample pseudo-symmetric neighbourhoods. A representation built around one canonical cell can make such transitions look like jumps in description. A lattice-level coordinate instead allows the physical trajectory and the changing set of exact basis correspondences to be treated separately.

\subsection{Exact transformations make the representation operational}

The invariant would be of limited use to crystallographic software if the path back to basis-dependent data were approximate or ambiguous. Retaining the conventional-to-primitive rational map and subsequent unimodular integer operations closes that loop. Once an exact basis operator is available, the same algebra can be propagated through the rest of the crystallographic data model. Miller indices transform contragrediently; fractional coordinates transform in direct space; symmetry operators can be conjugated; orientation and reciprocal-basis matrices can be updated; and atomic models or reflection arrays can remain attached to the lattice operation that produced them.

This is one reason to keep the comparison layer and the algebra layer distinct in software. The comparison layer can be compact, Euclidean and aggressively indexed. The algebra layer can be exact and explicit. The two have different numerical requirements and should not be forced into one representation.

\subsection{Scope, attribution and required closure validation}

Three mathematical objects must remain distinct. Kurlin's structured root invariant is used for the complete isometry classification. The sorted six-root-product key used here is continuous and Euclidean but many-to-one. The exact operator test over the finite Selling-superbase closure establishes the crystallographic identity used for reindexing. No claim in this paper requires one object to be simultaneously complete, globally Euclidean and the sole identity test.

The attribution follows the same separation. Andrews--Bernstein--Sauter supply the crystallographic $S^6$ embedding, the multiset uniqueness of the six Selling scalars, and the practical reduction algorithm~\cite{ab1988,abs2019,ab2023}. Kurlin's April 2026 manuscript provides the type-dependent enumeration of obtuse superbases, proves completeness of the 3D root invariant and bounds changes of individual root products; it explicitly postpones the construction of a metric on the 3D invariant space to subsequent work. The continuity at sorting ties and the lower-bound equation~\eqref{eq:lowerbound} are elementary properties of the sorted search key used here. We therefore cite ABS for the Selling/multiset prior art, and Kurlin for the classification and root-product results, not for the $L^2$ search metric.

The typed Selling-superbase closure has been checked against exhaustive obtuse-superbase enumeration and against coset sizes on high-symmetry self-matches. For a primitive orthorhombic cell the closure contains 32 superbases in four isometry classes; matching over the \emph{full} closure recovers the order-8 mmm coset on self-reindexing and the same eight operators when one copy is presented in an even-class basis, whereas matching only class representatives under-counts the coset (2 instead of 8). For the CXIDB~83 hexagonal cell the full closure recovers the order-24 $6/mmm{\cdot}(-I)$ coset. Exact archive metrics may set the near-zero expander to zero; measured serial frames use a small relative boundary tolerance (default $10^{-3}$) and, when needed, an angular-noise zero tolerance so that noisy $P6_3$ frames are classified as V4 rather than V1. The full-PDB timing benchmark must likewise be rerun with the exact one-sorted-key-per-lattice indexing policy reported explicitly.

An embedding of sorted-key distances is a visualization, not an independent validation that nearby lattices are nearby. Its validity has two separable questions. First, does the sorted-key quotient preserve the local neighbourhood structure of the observed population well enough for the intended descriptive use? This is answered by the closure-aware collision/edge audit described above. Second, does structure seen in that representation correspond to anything outside the metric used to construct it? Meaningful external tests include whether crystal-system labels occupy coherent strata without having been supplied to the embedding, whether independent experimental metadata explain directions of variation, whether controlled trajectories are retrieved without boundary-induced loss, whether exact basis operators are recovered across reduction and Voronoi-type boundaries, and whether closure-complete search remains practical at archive scale. The first tests representation fidelity on the sampled data; the second tests scientific interpretation. Neither requires the sorted key to be a globally complete invariant.

Finally, the framework addresses equality, similarity and basis equivalence of primitive lattices. Supercell/subcell relations change the lattice itself and therefore require explicit enumeration of finite-index sublattices or superlattices when those relations are part of the search problem. They should not be conflated with alternative primitive bases of one lattice.

\section{Conclusion}

A century of crystallographic reduction practice has understandably emphasized the question: \emph{which cell should represent this lattice?} That question remains essential for classification and reporting. Large-scale computation benefits from separating retrieval from certification.

Selling reduction supplies the finite basis-aware object: a type-dependent closure of obtuse superbases, not universally a single 24-member orbit. Kurlin's classification tells us what that closure must contain. For high-volume search we then intentionally forget some of its structure and use the sorted six root products as a continuous Euclidean key. The resulting distance is a lower bound to any allowed relabelling-orbit distance, so radius search is conservative. Collisions are expected and harmless because identity is never assigned by the key alone. Returned candidates are certified through the complete Selling closure and exact rational/integer basis algebra.

The resulting rule is concise:
\begin{quote}
\textbf{Search continuously and conservatively; decide exactly.}
\end{quote}
Or, computationally,
\begin{quote}
\textbf{Euclidean key $\rightarrow$ finite Selling closure $\rightarrow$ exact operator.}
\end{quote}
This architecture turns the expensive retrieval stage of lattice comparison into ordinary computational geometry without confusing a convenient search coordinate with a complete lattice metric, and it preserves the exact transformations needed by Miller indices, coordinates, symmetry and reindexing.

\section*{Data and code availability}
The implementation, CrystFEL stream parser, benchmark inputs and scripts generating the figures will be deposited in a versioned public repository before submission. The lysozyme cells derive from public PDB entries (UniProt P00698). The serial data are from a $\beta$-lactamase XFEL data set deposited as CXIDB entry 83. [Authors to confirm accession and add primary citation.] [Insert repository URL and archival DOI.]

\section*{Acknowledgements}
[Insert funding and acknowledgements.]

\begin{thebibliography}{99}
\small
\bibitem{ab1988} Andrews, L. C. \& Bernstein, H. J. (1988). \textit{Acta Cryst.} A\textbf{44}, 1009--1018.
\bibitem{ab2014} Andrews, L. C. \& Bernstein, H. J. (2014). \textit{J. Appl. Cryst.} \textbf{47}, 346--359.
\bibitem{ab2016} Andrews, L. C. \& Bernstein, H. J. (2016). \textit{J. Appl. Cryst.} \textbf{49}, 756--761.
\bibitem{abs2019} Andrews, L. C., Bernstein, H. J. \& Sauter, N. K. (2019). \textit{Acta Cryst.} A\textbf{75}, 593--599.
\bibitem{ab2023} Andrews, L. C. \& Bernstein, H. J. (2023). \textit{Acta Cryst.} A\textbf{79}, 485--498.
\bibitem{bravais1850} Bravais, A. (1850). M\'emoire sur les syst\`emes form\'es par des points distribu\'es r\'eguli\`erement sur un plan ou dans l'espace. \textit{J. \'{E}cole Polytechnique}, \textbf{19}, 1--128.
\bibitem{bd2014} Brehm, W. \& Diederichs, K. (2014). \textit{Acta Cryst.} D\textbf{70}, 101--109.
\bibitem{bck2021} Bright, M., Cooper, A. I. \& Kurlin, V. (2021). A complete and continuous map of the lattice isometry space for all 3-dimensional lattices. \texttt{arXiv:2109.11538}.
\bibitem{delone1932} Delone, B. N. (1932). Neue Darstellung der geometrischen Kristallographie. \textit{Z. Kristallogr.} \textbf{84}, 109--149.
\bibitem{fiedler1973} Fiedler, M. (1973). \textit{Czechoslovak Math. J.} \textbf{23}, 298--305.
\bibitem{gw2018} Gildea, R. J. \& Winter, G. (2018). \textit{Acta Cryst.} D\textbf{74}, 405--410.
\bibitem{kurlin2026} Kurlin, V. (2026). A complete isometry classification of 3-dimensional lattices. Revised manuscript, April 2026. \url{https://kurlin.org/projects/lattice-geometry/lattices3Dmaths.pdf}.
\bibitem{mcgill2014} McGill, K. J., Asadi, M., Karakasheva, M. T., Andrews, L. C. \& Bernstein, H. J. (2014). \textit{J. Appl. Cryst.} \textbf{47}, 360--364.
\end{thebibliography}

\end{document}
